\section{Introduction}

GNN inference and graph analytics on large graphs run partition-parallel. Each of $k$ workers
processes one part of the graph, and the workers exchange boundary values after every
superstep \cite{malewicz2010pregel}. Performance depends on load balance, on communication
volume, and on how cores are assigned to workers. METIS \cite{karypis1998metis} optimises the
partitioning for a single snapshot, but real graphs are \emph{dynamic}
\cite{dossantos2025landscape,zhang2026survey}. Edges and vertices arrive and expire, and the
active region moves, so partitions develop hotspots and a growing cut. A runtime can react by
shifting threads (cheap, but this fixes only imbalance), by refining the partitioning locally,
or by repartitioning from scratch (best quality, but it migrates much of the graph). Which
reaction pays off depends on whether a change is transient or persistent and on the cost of
communication. These factors are hard to capture with fixed rules.

We present \textbf{GNN-AdaptGraph}, an adaptive runtime in which a GraphSAGE controller
\cite{hamilton2017graphsage} reads the communication graph of the current partitioning and
decides, in every window, whether and how strongly to adapt. Compared with our earlier
version, the controller is trained on the \emph{outcomes} of actions instead of a hand-written
rule, the repartitioner optimises true communication volume, thread adjustment is
implemented and measured, and the evaluation uses real dynamic graphs and established
baselines. Our contributions are:
\begin{itemize}
  \item \textbf{A cost-explicit formulation of runtime adaptation} for partition-parallel
        processing of dynamic graphs (Sect.~\ref{sec:model}). The objective is end-to-end time,
        which includes adaptation and migration overheads, so no hand-tuned reward weights are needed.
  \item \textbf{A GraphSAGE controller trained on outcome-based oracle labels with a
        cost-sensitive (expected-regret) loss} (Sect.~\ref{sec:labels}). It chooses among four
        actions of increasing cost: \NOA{}, \TA{}, \RB{} and \RP{}.
  \item \textbf{ACGR}, a communication-aware local repartitioner with an explicit balance
        constraint and incremental $O(\deg v)$ move evaluation (Sect.~\ref{sec:acgr}).
  \item \textbf{A partition-parallel executor} in which \TA{} really changes per-partition
        parallelism, and a cost model calibrated on it (held-out $R^2=\CalRtwo$).
  \item \textbf{An evaluation} on a real temporal graph and on synthetic drift/hotspot
        streams, against METIS, KL, LDG, FENNEL, label propagation, periodic METIS
        repartitioning and a tuned threshold rule, on three platforms, with repeated
        measurements and significance tests (Sect.~\ref{sec:results}). It shows when each
        adaptation pays off and where learned control helps and where it does not.
\end{itemize}
